\documentclass[11pt,a4paper]{article}
\usepackage{fontspec}
\directlua{
  luaotfload.add_fallback("hbfb", {"DejaVu Serif:mode=node;", "DejaVu Sans:mode=node;"})
  luaotfload.add_fallback("ltfb", {"DejaVu Sans:mode=node;"})
}
\usepackage[bidi=basic]{babel}
\babelprovide[import, main]{hebrew}
\babelprovide[import]{english}
\babelfont{rm}[Ligatures=TeX, RawFeature={fallback=ltfb}]{DejaVu Serif}
\babelfont[hebrew]{rm}[Script=Hebrew, Ligatures=TeX, RawFeature={fallback=hbfb}]{Noto Serif Hebrew}
\babelfont{tt}[Scale=0.84, RawFeature={fallback=ltfb}]{DejaVu Sans Mono}

\usepackage[a4paper,top=2.1cm,bottom=2.1cm,left=1.9cm,right=1.9cm,headsep=0.5cm,footskip=1cm]{geometry}
\usepackage{booktabs, longtable, array, adjustbox, xcolor, enumitem, fancyhdr, titlesec, fancyvrb, needspace}

\definecolor{inkblue}{HTML}{14365A}
\definecolor{rulegray}{HTML}{8A97A3}
\definecolor{boxgray}{HTML}{F4F6F8}
\definecolor{codebg}{HTML}{F7F8FA}

\setlength{\parindent}{0pt}
\setlength{\parskip}{3pt plus 1pt}
\renewcommand{\arraystretch}{1.15}

% ---- running head / foot ----
\pagestyle{fancy}\fancyhf{}
\renewcommand{\headrulewidth}{0.4pt}
\renewcommand{\footrulewidth}{0pt}
\fancyhead[R]{\small בסיסי נתונים ו-SQL \textbar{} מודול 16}
\fancyhead[L]{\small דף עבודה מס' 1}
\fancyfoot[C]{\small עמוד \thepage{} מתוך \pageref{LastPage}}
\usepackage{lastpage}

% ---- section titles ----
\titleformat{\section}{\needspace{3\baselineskip}\large\bfseries\color{inkblue}}{}{0pt}{}
\titlespacing*{\section}{0pt}{14pt plus 3pt}{6pt}

\newcommand{\partheader}[3]{%
  \section*{#1}%
  \vspace{-4pt}{\color{rulegray}\hrule height 0.8pt}\vspace{4pt}%
  {\small\textbf{בסיס הנתונים:} #2 \hfill \textbf{ניקוד:} #3}\par\vspace{2pt}%
  {\color{rulegray}\hrule height 0.4pt}\vspace{6pt}}

% ---- a question ----
\newcounter{qnum}[section]
\newcommand{\question}[2]{%
  \needspace{5\baselineskip}\refstepcounter{qnum}\par\vspace{7pt}%
  \noindent\textbf{\color{inkblue}שאלה \theqnum.}\ \textbf{#2}\hfill{\small(#1 נק')}\par\vspace{2pt}}
\newcommand{\subq}[1]{\par\vspace{1pt}\hspace{1.1em}\textbf{#1.}~}

% ---- dotted answer space ----
\newcommand{\answerlines}[1]{%
  \par\vspace{1pt}%
  \begingroup\color{rulegray}%
  \count255=#1\relax
  \loop\ifnum\count255>0
    \noindent\babelsublr{\rule[-2pt]{0.6pt}{9pt}\rule[-2pt]{\dimexpr\linewidth-0.6pt}{0.3pt}}\par\vspace{5.4mm}%
    \advance\count255 by -1
  \repeat
  \endgroup\vspace{-2mm}}

% ---- SQL code, always left-to-right ----
\DefineVerbatimEnvironment{sqlverb}{Verbatim}{fontsize=\small, frame=leftline, framerule=1.2pt,
  rulecolor=\color{rulegray}, framesep=7pt, xleftmargin=12pt, samepage=true}
\newsavebox{\panelbox}
\newenvironment{panel}{\par\addvspace{7pt}\begin{lrbox}{\panelbox}%
  \begin{minipage}{\dimexpr\linewidth-2\fboxsep-2\fboxrule}}%
  {\end{minipage}\end{lrbox}\noindent\fcolorbox{rulegray}{boxgray}{\usebox{\panelbox}}\par\addvspace{7pt}}

% ---- data preview table ----
\newenvironment{tablepreview}[2]{%
  \needspace{4\baselineskip}\par\vspace{4pt}%
  \noindent{\small הטבלה \babelsublr{\texttt{\textbf{#1}}} \textendash\ #2 שורות}\par\vspace{1pt}%
  \begin{otherlanguage}{english}\begin{adjustbox}{max width=\linewidth}}
  {\end{adjustbox}\end{otherlanguage}\par\vspace{3pt}}
\newcommand{\thd}[1]{\textbf{#1}}
\newenvironment{fulltable}[3]{%
  \needspace{5\baselineskip}\par\addvspace{7pt}%
  \noindent{\small הטבלה \babelsublr{\texttt{\textbf{#1}}} \textendash\ #2 שורות#3}\par\vspace{2pt}%
  \begin{otherlanguage}{english}\begin{adjustbox}{max width=\linewidth}}
  {\end{adjustbox}\end{otherlanguage}\par\addvspace{4pt}}

\providecommand{\tightlist}{\setlength{\itemsep}{0pt}\setlength{\parskip}{0pt}}

\begin{document}


\begin{center}
{\LARGE\bfseries\color{inkblue} בסיסי נתונים ושפת SQL}\\[2pt]
{\large מודול 16 \textendash\ המשפט הראשון: \texttt{SELECT}, כינויים, \texttt{DISTINCT} ו-\texttt{NULL}}\\[6pt]
{\large\bfseries דף עבודה להגשה}
\end{center}
\vspace{2pt}
{\color{inkblue}\hrule height 1.2pt}
\vspace{6pt}

\noindent שם המגיש/ה:~\rule{5.2cm}{0.4pt}\hfill כיתה:~\rule{2.2cm}{0.4pt}\hfill תאריך:~\rule{3cm}{0.4pt}
\vspace{8pt}

\begin{panel}
\textbf{הוראות}
\begin{itemize}[leftmargin=1.2em, itemsep=1pt, topsep=3pt]
  \item בדף שני חלקים, ובכל חלק בסיס נתונים אחר. \textbf{סך הכול 100 נקודות.}
  \item לכל חלק מצורפות \textbf{הטבלאות ושלוש הרשומות הראשונות} שבכל טבלה. אלה הנתונים שעליהם נשאלות השאלות.
  \item את השאילתות כותבים בשורות שמתחת לכל שאלה. \textbf{שפת SQL נכתבת משמאל לימין} \textendash\ מתחילים מהסימן שבקצה \textbf{השמאלי} של השורה.
  \item לפני ההגשה \textbf{מריצים} כל שאילתה: \texttt{onecompiler.com/sqlite} או \texttt{sqliteonline.com}.
  \item כתבו שם טבלה ושם עמודה במלואם, בסגנון \texttt{Table.Column}, וסיימו כל שאילתה בנקודה-פסיק.
  \item ערך טקסט נכתב בין גרשיים \textbf{בודדים}: \texttt{'Haifa'}. גרשיים כפולים הם \textbf{שם} של עמודה או כינוי.
  \item בכל חישוב תאריכים השתמשו בתאריך הייחוס \texttt{'2026-09-21'}, ולא ב-\texttt{DATE('now')}.
\end{itemize}
\end{panel}


\vspace{10pt}
\section*{תזכורת לפני שמתחילים}
\vspace{-6pt}{\color{rulegray}\hrule height 0.8pt}\vspace{6pt}

\textbf{מבנה השאילתה.} שני חלקים, תמיד בסדר הזה, ונקודה-פסיק בסוף:

\begin{otherlanguage}{english}
\begin{sqlverb}
SELECT   <what to show>      -- the columns
FROM     <which table>;      -- the source
\end{sqlverb}
\end{otherlanguage}

\vspace{2pt}
\textbf{מה אפשר לכתוב ב-\babelsublr{\texttt{SELECT}}} (הדוגמאות על בסיס הנתונים של חלק א'):

\vspace{2pt}
\noindent\begin{tabular}{@{}p{0.69\linewidth}p{0.27\linewidth}@{}}
\toprule
\thd{איך זה נראה} & \thd{מה רוצים} \\
\midrule
\babelsublr{\texttt{SELECT * FROM Students;}} & כל העמודות \\
\babelsublr{\texttt{SELECT FirstName, LastName FROM Students;}} & עמודות נבחרות \\
\babelsublr{\texttt{SELECT Salary AS AnnualPay FROM Teachers;}} & כותרת אחרת לעמודה \\
\babelsublr{\texttt{SELECT Salary AS "שכר חודשי" FROM Teachers;}} & כותרת עם רווח או בעברית \\
\babelsublr{\texttt{SELECT DISTINCT Subject FROM Teachers;}} & בלי כפילויות \\
\babelsublr{\texttt{SELECT Salary * 12 AS Yearly FROM Teachers;}} & חישוב \\
\babelsublr{\texttt{SELECT FirstName || ' ' || LastName FROM Students;}} & חיבור מחרוזות \\
\bottomrule
\end{tabular}

\vspace{8pt}
\begin{panel}
\textbf{שלושה כללים שחוזרים בכל שאלה}
\begin{enumerate}[leftmargin=1.4em, itemsep=2pt, topsep=3pt]
  \item \textbf{גרש בודד לערך, גרש כפול לשם.} \babelsublr{\texttt{'Haifa'}} הוא טקסט; \babelsublr{\texttt{"Haifa"}} הוא \textbf{שם} של עמודה או של כינוי \textendash\ ולכן \babelsublr{\texttt{SELECT "Haifa" FROM Students}} ייכשל.
  \item \textbf{\babelsublr{\texttt{NULL}} מדביק.} כל פעולה עם \babelsublr{\texttt{NULL}} מחזירה \babelsublr{\texttt{NULL}}, גם חיבור מחרוזות. ערך חסר אינו אפס ואינו מחרוזת ריקה.
  \item \textbf{כל שאילתה מחזירה טבלה} \textendash\ והטבלה המקורית לא משתנה. \babelsublr{\texttt{SELECT}} רק מציג.
\end{enumerate}
\end{panel}

\needspace{9\baselineskip}\vspace{2pt}
\textbf{ארבע הטעויות שכדאי לא לעשות}

\vspace{2pt}
\noindent\begin{tabular}{@{}p{0.33\linewidth}p{0.30\linewidth}p{0.30\linewidth}@{}}
\toprule
\thd{התיקון} & \thd{מה קורה} & \thd{הטעות} \\
\midrule
גרש בודד & מחפש עמודה בשם הזה & גרש כפול למחרוזת \\
להסיר את הפסיק & \babelsublr{\texttt{syntax error near FROM}} & פסיק לפני \babelsublr{\texttt{FROM}} \\
\babelsublr{\texttt{|| ' ' ||}} באמצע & \babelsublr{\texttt{AdamHalabi}} & חיבור בלי רווח \\
לזהות, ולטפל במודול 20 & כל הביטוי נעלם & לצפות ש-\babelsublr{\texttt{NULL}} יתנהג כמו 0 \\
\bottomrule
\end{tabular}

\vspace{6pt}
{\small\textbf{נתקעתם בשגיאה?} קראו את המילה שאחרי \babelsublr{\texttt{near}} \textendash\ שם נמצאת הבעיה, או ממש לפניה.}

\vspace{6pt}
\partheader{חלק א' \textendash\ בית הספר "עתיד"}{\texttt{school.sql}}{84 נקודות (12 שאלות \texttt{$\times$} 7)}
\textbf{הנתונים.} שלוש הרשומות הראשונות בכל טבלה:\par\vspace{2pt}
\begin{tablepreview}{Cities}{6}
\begin{tabular}{ll}
\toprule
\thd{CityCode} & \thd{CityName} \\
\midrule
1 & ירכא \\
2 & עכו \\
3 & כרמיאל \\
\bottomrule
\end{tabular}
\end{tablepreview}

\begin{tablepreview}{Teachers}{7}
\begin{tabular}{llllllll}
\toprule
\thd{TeacherCode} & \thd{FirstName} & \thd{LastName} & \thd{Subject} & \thd{HireDate} & \thd{Salary} & \thd{CityCode} & \thd{Phone} \\
\midrule
1 & נביל & סרחאן & מתמטיקה & 2012-09-01 & 14500 & 1 & 050-7010101 \\
2 & רונית & בר-לב & אנגלית & 2018-09-01 & 11200 & 4 & \textit{NULL} \\
3 & חוסאם & זיאד & מחשבים & 2021-02-15 & 12800 & 6 & 054-7030303 \\
\bottomrule
\end{tabular}
\end{tablepreview}

\begin{tablepreview}{Classes}{5}
\begin{tabular}{lllll}
\toprule
\thd{ClassCode} & \thd{ClassName} & \thd{Grade} & \thd{TeacherCode} & \thd{RoomNumber} \\
\midrule
101 & י1 & 10 & 1 & 12 \\
102 & י2 & 10 & 2 & 14 \\
103 & יא1 & 11 & 3 & 21 \\
\bottomrule
\end{tabular}
\end{tablepreview}

\begin{tablepreview}{Courses}{7}
\begin{tabular}{llll}
\toprule
\thd{CourseCode} & \thd{CourseName} & \thd{TeacherCode} & \thd{WeeklyHours} \\
\midrule
11 & מתמטיקה 5 יח"ל & 1 & 5 \\
12 & אנגלית 4 יח"ל & 2 & 4 \\
13 & מבוא לבסיסי נתונים & 3 & 3 \\
\bottomrule
\end{tabular}
\end{tablepreview}

\begin{tablepreview}{Students}{18}
\begin{tabular}{lllllllll}
\toprule
\thd{StudentId} & \thd{FirstName} & \thd{LastName} & \thd{ClassCode} & \thd{Gender} & \thd{BirthDate} & \thd{CityCode} & \thd{EnrollDate} & \thd{Phone} \\
\midrule
1001 & אדם & חלבי & 101 & M & 2011-03-14 & 1 & 2026-09-01 & 050-1000001 \\
1002 & נור & עזאם & 101 & F & 2011-07-22 & 1 & 2026-09-01 & \textit{NULL} \\
1003 & יואב & כהן & 101 & M & 2011-01-09 & 3 & 2026-09-01 & 052-1000003 \\
\bottomrule
\end{tabular}
\end{tablepreview}

\begin{tablepreview}{Grades}{60}
\begin{tabular}{llll}
\toprule
\thd{StudentId} & \thd{CourseCode} & \thd{Term} & \thd{Grade} \\
\midrule
1001 & 11 & 1 & 88 \\
1001 & 12 & 1 & 74 \\
1001 & 14 & 1 & 92 \\
\bottomrule
\end{tabular}
\end{tablepreview}

\begin{tablepreview}{Absences}{16}
\begin{tabular}{lllll}
\toprule
\thd{AbsenceId} & \thd{StudentId} & \thd{AbsenceDate} & \thd{Excused} & \thd{Reason} \\
\midrule
1 & 1003 & 2026-09-02 & 1 & מחלה \\
2 & 1003 & 2026-09-03 & 1 & מחלה \\
3 & 1004 & 2026-09-07 & 0 & \textit{NULL} \\
\bottomrule
\end{tabular}
\end{tablepreview}

\vspace{6pt}{\color{rulegray}\hrule height 0.4pt}
\setcounter{qnum}{0}
\question{7}{}
הציגו \textbf{את כל} העמודות של כל המורים. כמה שורות חזרו?
\answerlines{3}
\question{7}{}
הציגו \textbf{שם פרטי ושם משפחה} של כל התלמידים.
\answerlines{3}
\question{7}{}
הציגו את שם המקצוע ושעות הלימוד השבועיות, עם כינויים בעברית: \babelsublr{\texttt{מקצוע}} ו-\babelsublr{\texttt{שעות\ שבועיות}}.
\answerlines{3}
\question{7}{}
אילו \textbf{מקצועות הוראה} (\babelsublr{\texttt{Subject}}) קיימים בבית הספר? כל מקצוע \textbf{פעם אחת}. יש 7 מורים \textemdash{} למה התוצאה קטנה יותר?
\answerlines{3}
\question{7}{}
אילו \textbf{שכבות} (\babelsublr{\texttt{Grade}}) יש בבית הספר? בלי כפילויות.
\answerlines{3}
\question{7}{}
הציגו לכל מורה \textbf{שם מלא} בעמודה אחת (שם פרטי, רווח, שם משפחה), בכינוי \babelsublr{\texttt{FullName}}.
\answerlines{3}
\question{7}{}
הציגו לכל מורה את שם המשפחה, השכר החודשי, ו\textbf{השכר השנתי} המחושב (×12) בכינוי \babelsublr{\texttt{AnnualSalary}}.
\answerlines{3}
\question{7}{}
שנת לימודים היא 30 שבועות. הציגו לכל מקצוע \textbf{כמה שעות לימוד בשנה}, בכינוי \babelsublr{\texttt{YearlyHours}}.
\answerlines{3}
\question{7}{}
הציגו שם פרטי וטלפון של כל התלמידים. \textbf{אצל כמה מהם} הטלפון הוא \babelsublr{\texttt{NULL}}? (ספרו בעיניים \textemdash{} \babelsublr{\texttt{COUNT}} נלמד בשיעור 23.)
\answerlines{3}
\question{7}{}
\textbf{מלכודת.} הריצו:

\begin{otherlanguage}{english}
\begin{sqlverb}
SELECT Students.FirstName, Students.Phone || ' (בית)' AS PhoneLabel FROM Students;
\end{sqlverb}
\end{otherlanguage}

מה קיבלתם אצל התלמידים שאין להם טלפון? \textbf{הסבירו למה} \textemdash{} ולא, זו לא שגיאה.
\answerlines{3}
\question{7}{}
הציגו את \textbf{5 השורות הראשונות} של טבלת הציונים. מה אתם מסיקים על \textbf{המפתח הראשי} שלה?
\answerlines{3}
\question{7}{}
\textbf{קריאת שגיאה.} הריצו:

\begin{otherlanguage}{english}
\begin{sqlverb}
SELECT Students.FirstName, Students.LastName FROM Student;
\end{sqlverb}
\end{otherlanguage}

מה ההודעה? מה המילה החשובה בה, ומה התיקון?
\answerlines{3}

\clearpage
\partheader{חלק ב' \textendash\ חוגי ספורט \textendash\ בסגנון שאלון 735911}{\texttt{clubs.sql}}{16 נקודות}
\textbf{הנתונים.} שלוש הרשומות הראשונות בכל טבלה:\par\vspace{2pt}
\begin{tablepreview}{Cities}{6}
\begin{tabular}{ll}
\toprule
\thd{CityCode} & \thd{CityName} \\
\midrule
1 & ירכא \\
2 & עכו \\
3 & כרמיאל \\
\bottomrule
\end{tabular}
\end{tablepreview}

\begin{tablepreview}{Coaches}{5}
\begin{tabular}{llll}
\toprule
\thd{CoachCode} & \thd{FirstName} & \thd{LastName} & \thd{Phone} \\
\midrule
1 & סאמר & חלבי & 050-1111111 \\
2 & מיכל & לוי & \textit{NULL} \\
3 & ראמי & עזאם & 054-3333333 \\
\bottomrule
\end{tabular}
\end{tablepreview}

\begin{tablepreview}{Clubs}{8}
\begin{tabular}{lllll}
\toprule
\thd{ClubCode} & \thd{ClubName} & \thd{Sport} & \thd{CoachCode} & \thd{Price} \\
\midrule
100 & כדורגל צעירים & כדורגל & 1 & 150 \\
101 & כדורגל בוגרים & כדורגל & 1 & 180 \\
102 & שחייה מתחילים & שחייה & 2 & 220 \\
\bottomrule
\end{tabular}
\end{tablepreview}

\begin{tablepreview}{Members}{12}
\begin{tabular}{llllll}
\toprule
\thd{Id} & \thd{FirstName} & \thd{LastName} & \thd{CityCode} & \thd{Age} & \thd{JoinYear} \\
\midrule
1001 & אדם & חלבי & 1 & 14 & 2023 \\
1002 & נור & עזאם & 1 & 15 & 2024 \\
1003 & יואב & כהן & 3 & 16 & 2022 \\
\bottomrule
\end{tabular}
\end{tablepreview}

\begin{tablepreview}{ClubMembers}{18}
\begin{tabular}{ll}
\toprule
\thd{Id} & \thd{ClubCode} \\
\midrule
1001 & 100 \\
1001 & 104 \\
1002 & 102 \\
\bottomrule
\end{tabular}
\end{tablepreview}

\vspace{6pt}{\color{rulegray}\hrule height 0.4pt}
\setcounter{qnum}{0}
\question{4}{}
איזו פקודה מחזירה את \textbf{כל השדות} של כל הרשומות בטבלה Clubs?

\begin{enumerate}
\def\labelenumi{\arabic{enumi}.}
\tightlist
\item
  \babelsublr{\texttt{SELECT\ All\ FROM\ Clubs}}
\item
  \babelsublr{\texttt{SELECT\ *\ FROM\ Clubs}}
\item
  \babelsublr{\texttt{SELECT\ Clubs\ FROM\ *}}
\item
  \babelsublr{\texttt{SELECT\ (ALL)\ FROM\ Clubs}}
\end{enumerate}
\answerlines{3}
\question{6}{}
כמה שורות תחזיר השאילתה \babelsublr{\texttt{SELECT\ DISTINCT\ Clubs.Sport\ FROM\ Clubs;}}? ומה הייתה התשובה בלי \babelsublr{\texttt{DISTINCT}}?
\answerlines{3}
\question{6}{}
כתבו בסגנון הבחינה שאילתה שמציגה את שם החוג ואת המחיר שלו, כשהכותרת של עמודת המחיר היא "מחיר חודשי".
\answerlines{3}

\end{document}